% --------------------------------------------------------------------------
% How to Use This Manual.
% --------------------------------------------------------------------------
\chapter*{How to Use This Manual}
\addcontentsline{toc}{chapter}{How to Use This Manual}
\markboth{How to Use This Manual}{How to Use This Manual}

This manual is in two parts and a set of appendices.

\textbf{Part One, \emph{Using \ccprod},} is written to be read. It starts at
the screen and works outwards: what is on the display, how the keyboard moves
the cursor, how a value gets into a cell, what the menus contain, how files are
opened and saved, and what the commands that act on more than one cell will and
will not do. Read it once, in order, and you will not need to read it again.

\textbf{Part Two, \emph{Reference},} is written to be looked things up in.
Every menu command, every key, every function and every message the program can
display has an entry, and each entry stands on its own. Nothing in Part Two
assumes you have read Part One, and nothing in it is a tutorial.

\textbf{The appendices} are the numbers: how much the program can hold, what a
workbook file contains byte for byte, how a number is represented, and where
the machine's memory goes. You do not need any of it to use \ccprod. You need
all of it to know what the program will do at the edges.

\section*{What this manual assumes}

That you have used a spreadsheet before --- that a cell, a row, a column and a
formula are already familiar. \ccprod{} is small, and this manual is
correspondingly short on explanations of what a spreadsheet is for.

It does \emph{not} assume that the spreadsheet you have used before was this
one. \ccprod{} is not a clone of anything. It leaves out a great deal that
other programs have, and a few of the things it does keep, it does differently.
Where that is true, this manual says so plainly rather than leaving you to find
out. Appendix~\ref{app:differences} collects every such difference in one list.

\section*{Conventions}

\vspace{-4pt}
\begin{cctable}{Convention}{Meaning}{1.42in}{3.4in}
\key{F7} & A key on the keyboard. Press it. \\
\keyc{Ctrl}{C} & Two keys pressed together: hold the first, press the second. \\
\menupath{File > Import} & A menu path. Open the \menupath{File} menu, then
choose \menupath{Import}. \\
\lit{=SUM(A1:A10)} & Something you type, or something the program displays,
character for character. \\
\msg{no matching files on the card} & A message the program puts on the screen.
Messages are quoted exactly, including their punctuation and their spelling. \\
\errv{\#DIV/0!} & An error value, as it appears in a cell. \\
\fname{CMDRCALC.PRG} & A file name. \\
\cellref{A1}, \cellref{B2:D40} & A cell and a range of cells. \\
\end{cctable}

Two kinds of aside interrupt the text.

\begin{ccnote}
A note tells you something useful that is not obvious, usually because it is a
consequence of how the program is built rather than of how it is used.
\end{ccnote}

\begin{cccaution}
A caution tells you about something that costs you work: a command that cannot
be undone, a limit that silently truncates, or a file that will not come back.
Cautions are rare, and every one of them is there because the alternative is
losing a workbook.
\end{cccaution}

\section*{Keys and the Commander X16 keyboard}

\ccprod{} puts the machine into ISO mode at startup, which is what makes the
keyboard produce ordinary ASCII: unshifted letters arrive lowercase, shifted
letters uppercase. Everything in this manual that talks about typing assumes
that.

Several keys on the Commander X16 keyboard send the same code as a control
combination, and the program cannot tell them apart. \keyc{Ctrl}{M} is
\key{Return}; \keyc{Ctrl}{I} is \key{Tab}. The full list is in Chapter~\ref{ch:keyref}.
Where this manual names a key, it names the one you are most likely to be
pressing.

\key{F12} is not used by \ccprod{} and never will be: on the emulator it opens
the machine's own debugger, and a program that took the key would take that
away.

\section*{The page numbers}

Pages are numbered by chapter. Page \lit{9-4} is the fourth page of Chapter~\ref{ch:importexport};
page \lit{B-2} is the second page of Appendix~\ref{app:fileformat}. The black tab on the outer
edge of each page steps one notch down the page per chapter, so the manual can
be opened at a chapter by fanning it rather than by reading the numbers.

\clearpage
